\item All names that are referenced in a program need to be declared
  regardless of whether the reference gets evaluated or not.
\begin{lstlisting}
print(1 if True else x)
\end{lstlisting}
Leads to an error \texttt{name not found in the current or enclosing environment(s).}

A name is declared by an assignment, a \texttt{for} loop, a parameter or a function definition.
The scope of a name that is declared in a function is that function, including the functions
nested in it; a name that is declared at the top level is in scope everywhere in the program.
A name that is declared in one function can therefore not be used in another function, unless
the second function is nested in the first. A \texttt{global} declaration, anywhere in the
program, declares the name at the top level (and not in the function that contains it), so that
every function can use it, also a function that comes before the one with the declaration.
Built-in names are in scope everywhere. The order of the declarations and the uses does not
matter: a name that is used before it is assigned is not rejected, but leads to a run-time
\texttt{NameError}, as in Python.
\begin{lstlisting}
def swap(y):
    print(x)    # x is declared global in g: swap can be defined first

def g():
    global x
\end{lstlisting}

\item In Python, a logical line is the smallest unit of code that the 
interpreter can execute, typically corresponding to a complete statement. 
A physical line, by contrast, is a single line in the source file, 
usually ending with a newline character. A logical line may span multiple 
physical lines depending on whether line joining is used.

By default, each physical line corresponds to a logical line. However, 
when a statement is too long, Python allows it to be split across multiple 
physical lines using line continuation. Python provides two ways to join 
lines: \textbf{explicit line joining} and \textbf{implicit line joining}.

\textbf{Explicit line joining} uses a backslash (\textbackslash) to join 
the current physical line with the next one, forming a single logical line. 
The backslash must be the last character on the line—no characters 
(including whitespace or comments) are allowed after it.

\begin{lstlisting}
total = 1 + 2 + 3 + \
        4 + 5
\end{lstlisting}

\textbf{Implicit line joining} occurs when expressions are enclosed in 
parentheses \texttt{()}, square brackets \texttt{[]}, or braces \texttt{\{\}}. 
Within these delimiters, line breaks are allowed without using a backslash.

\begin{lstlisting}
total = (
            1 + 2 + 3 +
            4 + 5
         )
\end{lstlisting}

Therefore, when breaking a statement across lines, you must either use a backslash for explicit continuation or wrap the expression in parentheses, brackets, or braces for implicit continuation. Python requires each line to be syntactically complete—it does not infer the end of a statement implicitly.

For details on logical lines, see the official Python documentation:
\href{https://docs.python.org/3/reference/lexical_analysis.html#logical-lines}{Python Lexical Analysis – Logical Lines}.

\item Return statements are only allowed in bodies of functions. Each return statement must appear on a single logical line.
\item Lambda expressions are limited to a single logical line.
